%! Factoring Trinomials — Unit 1, Lesson 1.4 — Slides (Beamer)
\documentclass[aspectratio=169, 11pt]{beamer}
\usepackage{atda-beamer}   % provides \royalheader{} and \sectionlabel[color]{}

\begin{document}

% TITLE SLIDE (CourseName is not defined in beamer, so the name is literal)
{
\setbeamercolor{background canvas}{bg=royal}
\begin{frame}[plain]
  \vspace{2.8cm}\hspace{0.55cm}%
  \begin{minipage}{0.9\textwidth}
    {\color{goldacc}\bfseries\small LESSON 1.4}\\[0.5cm]
    {\color{white}\bfseries\LARGE Factoring Trinomials}\\[1.2cm]
    {\color{mistmid}\small Algebra, Trigonometry, and Data Analysis~$\cdot$~Unit 1}
  \end{minipage}
\end{frame}
}

% ── HOOK ──────────────────────────────────────────────────────────────────────
\begin{frame}[t, plain]
  \royalheader{The Banner with No Dimensions on the Order}
  \sectionlabel[goldacc]{hook --- how many guesses are you willing to make?}
  A sign shop's order lists only the area of a rectangular banner:
  \[ A(x) = 6x^2+19x+10 \quad \text{square feet.} \]
  Yesterday we recovered a patio's sides because the area had a shared factor. Check this one.
  \begin{block}{It has no common factor --- so 1.3's method stalls}
    The $6$ could split $1 \cdot 6$ or $2 \cdot 3$; the $10$ could split $1 \cdot 10$ or
    $2 \cdot 5$. How many guesses before you want a rule?
  \end{block}
\end{frame}

% ── WHERE THE RULE COMES FROM ─────────────────────────────────────────────────
\begin{frame}[t, plain]
  \royalheader{The Rule Is Hiding in the Multiplication}
  \sectionlabel{distribute first, then read it backwards}
  \[ (x+p)(x+q) = x^2 + \underbrace{(p+q)}_{\text{the middle}}x + \underbrace{pq}_{\text{the end}} \]
  To factor $x^2+bx+c$: find two integers whose \textbf{product} is $c$ and whose \textbf{sum}
  is $b$.
  \begin{block}{Signs, decided by $c$}
    $c>0$: both numbers carry the sign of $b$. \quad $c<0$: opposite signs, and the bigger one
    carries the sign of $b$.
  \end{block}
\end{frame}

% ── a = 1 PRACTICE ────────────────────────────────────────────────────────────
\begin{frame}[t, plain]
  \royalheader{Leading Coefficient $1$ --- All Four Sign Cases}
  \sectionlabel{product $c$, sum $b$}
  \vspace{0.2cm}
  \renewcommand{\arraystretch}{1.5}
  \begin{tabularx}{\textwidth}{@{}p{3.2cm} p{4.0cm} X@{}}
    \textbf{Trinomial} & \textbf{Two numbers} & \textbf{Factored} \\
    \hline
    $x^2+9x+20$  & $4$ and $5$    & $(x+4)(x+5)$ \\
    $x^2-7x+12$  & $-3$ and $-4$  & $(x-3)(x-4)$ \\
    $x^2+3x-10$  & $5$ and $-2$   & $(x+5)(x-2)$ \\
    $x^2-2x-15$  & $-5$ and $3$   & $(x-5)(x+3)$ \\
  \end{tabularx}
  \vspace{0.3cm}

  The product tells you the \textbf{signs}. The sum tells you \textbf{which pair}.
\end{frame}

% ── PRIME ─────────────────────────────────────────────────────────────────────
\begin{frame}[t, plain]
  \royalheader{When No Pair Works: \emph{Prime}}
  \sectionlabel[goldacc]{running out of options is a proof}
  \[ x^2+5x+7 \]
  There are exactly two integer pairs whose product is $7$:
  \[ 1 \cdot 7 = 7 \;\text{(sum } 8)  \qquad (-1)(-7) = 7 \;\text{(sum } {-8}) \]
  \begin{block}{Why this is certainty, not surrender}
    The list of factor pairs is \textbf{finite}. Checking all of them settles the question ---
    $x^2+5x+7$ is \textbf{prime} over the integers.
  \end{block}
\end{frame}

% ── ac METHOD ─────────────────────────────────────────────────────────────────
\begin{frame}[t, plain]
  \royalheader{When the Leading Coefficient Is Not $1$}
  \sectionlabel{your warm-up already found this rule}
  Three terms will not group. \textbf{Split the middle term} and you have four.
  \begin{enumerate}
    \setlength{\itemsep}{0.15cm}
    \item \textbf{GCF first.} Always --- it shrinks $ac$ too.
    \item Multiply $a \cdot c$.
    \item Find two integers with product $ac$ and sum $b$.
    \item Split $bx$ into those two terms.
    \item \textbf{Group} --- Lesson 1.3, unedited.
  \end{enumerate}
  \vspace{0.2cm}
  {\color{goldacc}\bfseries Only step 3 is new today.}
\end{frame}

% ── HOOK PAYOFF ───────────────────────────────────────────────────────────────
\begin{frame}[t, plain]
  \royalheader{Back to the Banner}
  \sectionlabel{$ac = 6 \cdot 10 = 60$, and $4+15 = 19$}
  \[
  \begin{aligned}
    6x^2+19x+10 &= 6x^2+4x+15x+10 \\
                &= 2x(3x+2) + 5(3x+2) \\
                &= (3x+2)(2x+5)
  \end{aligned}
  \]
  A banner $(3x+2)$ feet by $(2x+5)$ feet. At $x=2$: $8$ ft by $9$ ft, and $8 \cdot 9 = 72$ square
  feet.
  \begin{block}{The check costs ten seconds}
    Multiply it back out. If you do not land on $6x^2+19x+10$, you are not done.
  \end{block}
\end{frame}

% ── SIGNS ─────────────────────────────────────────────────────────────────────
\begin{frame}[t, plain]
  \royalheader{The Mixed-Sign Case}
  \sectionlabel[goldacc]{where the $ac$ method actually breaks}
  \[
  \begin{aligned}
    4x^2+4x-15 &= 4x^2+10x \;{\color{goldacc}-\;6x-15} \\
               &= 2x(2x+5) \;{\color{goldacc}-\;3(2x+5)} \\
               &= (2x+5)(2x-3)
  \end{aligned}
  \]
  Here $ac = -60$ and the pair is $10$ and $-6$. The negative comes out \textbf{front} of the
  second pair.
  \begin{block}{Diagnosis}
    If your two leftovers do not match, check this first --- writing $-3(2x-5)$ instead of
    $-3(2x+5)$ is nine failures out of ten.
  \end{block}
\end{frame}

% ── FACTOR COMPLETELY ─────────────────────────────────────────────────────────
\begin{frame}[t, plain]
  \royalheader{Factoring Completely: The Order, Updated}
  \sectionlabel{GCF first --- still}
  \begin{enumerate}
    \setlength{\itemsep}{0.2cm}
    \item \textbf{GCF first.} It makes every number after it smaller --- including $ac$.
    \item \textbf{Count the terms.} Four $\Rightarrow$ group. \textbf{Three $\Rightarrow$ today.}
          Two comes in 1.5.
    \item \textbf{Check every factor}, then multiply back.
  \end{enumerate}
  \vspace{0.2cm}
  \[
  \begin{aligned}
    3x^3+21x^2+30x &= 3x\left(x^2+7x+10\right) \;=\; 3x(x+5)(x+2)
  \end{aligned}
  \]
  That inner trinomial was last night's item 13.
\end{frame}

% ── ACTIVITY LAUNCH ───────────────────────────────────────────────────────────
\begin{frame}[t, plain]
  \royalheader{Group Activity: The Cafeteria Mural}
  \sectionlabel{groups of three --- everyone starts at Tier R}
  The district approved the mural by \textbf{area only} --- no drawing, no dimensions --- and the
  art club has to order border tape:
  \[ A(x) = 6x^2+17x+5 \quad \text{square feet.} \]
  \begin{block}{Your job}
    Recover the shape from the polynomial. The mural stands $(3x+1)$ feet tall --- which factor is
    the \textbf{width}, and how do you know?
  \end{block}
\end{frame}

% ── CLOSE ─────────────────────────────────────────────────────────────────────
\begin{frame}[t, plain]
  \royalheader{Exit Ticket}
  \sectionlabel[goldacc]{notes closed --- 5 minutes}
  \begin{enumerate}
    \setlength{\itemsep}{0.35cm}
    \item Factor: $x^2+9x+18$.
    \item Factor: $3x^2+11x+6$.
    \item A student factors $2x^2+16x+24$ as $(2x+4)(x+6)$. Correct? Complete? Explain.
  \end{enumerate}
  \vspace{0.3cm}
  {\color{royal}\bfseries Next: Lesson 1.5 --- Special Forms.}
\end{frame}

\end{document}
